\clearpage
\section{From joint observations to a comparison vector}\label{sec:representation}
\subsection{Follow the shapes}
One deterministically prepared example has shape
\[
 X\in\R^{2\times100\times17\times3}.
\]
The axes are person slot, sampled frame, joint, and $(x,y,\text{confidence})$.
A missing second person is represented by a padded slot. Temporal sampling
selects 100 frame positions; it is not a claim that the original recording
lasted 100 frames or that all clips had the same physical duration.

The preparation sequence is camera normalization, person tracking, conversion
from the COCO joint convention to the H36M joint convention, uniform temporal
sampling, padding, and spatial normalization. Tracking tries to keep person
slots consistent over time. Conversion produces the joint arrangement expected
by the pretrained model. Normalization removes some dependence on image position
and scale; it does not eliminate viewpoint, tracking, or pose-estimation errors.

Some converted joints combine several source joints. Their coordinates can be
averaged, but the local implementation uses the \emph{minimum} confidence of
the contributing joints. If one necessary source is absent, its derived joint
should not acquire high confidence merely because another source is visible.
This is an explicit departure from upstream confidence handling.

\begin{center}
\begin{tikzpicture}[node distance=6mm,
 box/.style={draw=ink,align=center,text width=113mm,minimum height=10mm,font=\small},
 arrow/.style={-{Latex},thick,ink}]
\node[box] (x) {Prepared poses: $2\times100\times17\times3$};
\node[box,below=of x] (f) {Frozen MotionBERT: $2\times100\times17\times512$};
\node[box,below=of f] (h) {Time mean, joint flattening, person mean: $h\in\R^{8704}$};
\node[box,below=of h] (u) {Learned affine head: $u=Wh+b\in\R^{2048}$};
\node[box,below=of u] (z) {Unit-length embedding: $z=u/\lVert u\rVert_2$};
\draw[arrow] (x)--(f);\draw[arrow] (f)--(h);\draw[arrow] (h)--(u);\draw[arrow] (u)--(z);
\end{tikzpicture}
\end{center}

\subsection{What MotionBERT contributes}
A representation is a numerical description useful for a later task. Raw
coordinates describe where joints were observed. Learned features can combine
positions across joints and time so that the next stage need not learn all
motion structure from scratch.

MotionBERT was pretrained to recover motion structure from partial or noisy 2D
observations using a 3D reconstruction objective. This study reuses its generic
pretrained encoder. The adapter requests its internal representation and bypasses
its 3D output head. There is no newly estimated 3D skeleton that must be evaluated
between our input poses and retrieval vectors.\cite{motionbert}

The essential transformer operation is attention. Put a set of token features
in a matrix $H\in\R^{L\times d}$, one token per row. Learned matrices produce
\[
 Q=HW_Q,\qquad K=HW_K,\qquad V=HW_V.
\]
The query--key dot product measures which tokens should exchange information.
For each row, softmax turns scores into positive weights that sum to one:
\[
 a_{ij}=\frac{\exp(q_i^\top k_j/\sqrt{d_k})}
 {\sum_\ell\exp(q_i^\top k_\ell/\sqrt{d_k})},
 \qquad o_i=\sum_j a_{ij}v_j.
\]
The square-root factor moderates score scale as the key dimension grows.
Unlike a fixed temporal average, these weights depend on the input. Multiple
attention heads provide several learned ways to mix information, while
position information distinguishes joints and time indices.

MotionBERT's spatial attention exchanges information among joints in a frame;
temporal attention exchanges information across frames. Its two streams apply
spatial and temporal operations in opposite orders and fuse their outputs.
You need to understand this role, not reproduce every transformer layer to
explain this experiment. The pretrained implementation is imported intact.

\textbf{Frozen} and \textbf{evaluation mode} are different controls. Frozen
parameters receive no training updates. Evaluation mode selects inference
behavior for modules with training-dependent behavior, such as dropout. The
project uses both, so all recipes share a fixed feature extractor.

\subsection{Why averaging features does not discard all motion information}
Let $F_{p,t,j,:}(X)\in\R^{512}$ denote an encoder output. The pre-head vector is
\[
 h(X)=\frac12\sum_{p=1}^{2}
       \operatorname{vec}_{j}\!\left(\frac1{100}
       \sum_{t=1}^{100}F_{p,t,j,:}(X)\right)\in\R^{17\cdot512}.
\]
Here $\operatorname{vec}_{j}$ concatenates the seventeen joint feature vectors.
The mean is taken after the encoder has combined information across time. It
is therefore not equivalent to averaging raw poses before any motion processing.
It still imposes an information bottleneck: whatever the pooled representation
does not retain cannot be recovered by the head.

With a frozen encoder, $h(X)$ can be calculated once and stored in a feature
cache. This saves repeated encoder work across head-training runs. A cache is
valid only for its exact source data, preprocessing, encoder weights, code,
and relevant runtime. A damaged query needs a fresh encoder pass on the damaged
tensor; adding noise to a cached vector would be a different intervention.

\subsection{The trainable head and the geometry of retrieval}
The trained transformation is
\[
 u=Wh+b,\qquad W\in\R^{2048\times8704},\quad b\in\R^{2048},
 \qquad z=\frac{u}{\lVert u\rVert_2}.
\]
There are $2048(8704+1)=17{,}827{,}840$ trainable parameters. The architecture
is simple, but its parameter count is not tiny. The cached $h$ vectors are
\emph{not} unit-normalized; the normalization follows the learned projection.
Zero projections are an invalid numerical state to detect rather than meaningful
directions for retrieval.

For nonzero vectors $u,v$, cosine similarity is
\[
 \cos(u,v)=\frac{u^\top v}{\lVert u\rVert\lVert v\rVert}.
\]
It compares direction rather than length. Once both vectors have unit norm,
$s_{ij}=z_i^\top z_j$. Expanding a squared distance gives
\[
 \lVert z_i-z_j\rVert^2
 =z_i^\top z_i+z_j^\top z_j-2z_i^\top z_j=2-2s_{ij}.
\]
Thus sorting by largest cosine is the same as sorting by smallest squared
Euclidean distance on these unit vectors. This equivalence is why the contextual
neighborhood can be expressed using $\max(2-2s_{ij},0)$. The maximum only guards
against a tiny negative distance from floating-point arithmetic.

\begin{concept}{A two-coordinate retrieval example}
Let a query be $q=(0.8,0.6)$ and the two gallery vectors be $g_1=(1,0)$ and
$g_2=(0,1)$. All are unit vectors. Their cosine scores are $0.8$ and $0.6$,
so $g_1$ ranks first. The squared distances are $0.4$ and $0.8$, producing
the same ranking. Multiplying the unnormalized query by ten changes neither
its normalized direction nor the ranking.
\end{concept}

The output has 2,048 coordinates, not 120 class probabilities. Labels influence
the geometry during learning; an unfamiliar action can still be recognized by
comparison with its gallery exemplar. Contextual neighborhoods affect training,
but \textbf{all three recipes use the same cosine retrieval at evaluation}.
There is no contextual reranking added at query time.

\subsection{From a loss to a weight update}
A loss $L(\theta)$ assigns a penalty to the current behavior, with
$\theta=(W,b)$. Minimizing it requires a direction in parameter space.
For a small perturbation $\delta$,
\[
 L(\theta+\delta)\approx L(\theta)+\nabla L(\theta)^\top\delta.
\]
Taking $\delta=-\eta\nabla L$ lowers this local approximation when $\eta>0$.
That is the basic gradient-descent idea. It does not promise global improvement
for an arbitrarily large step, or a global optimum for this nonconvex objective.

Backpropagation is the chain rule applied through the computation. For
$z=u/\lVert u\rVert$ away from zero,
\[
 \frac{\partial z}{\partial u}
 =\frac{I-zz^\top}{\lVert u\rVert}.
\]
The factor $I-zz^\top$ removes the component parallel to $z$: changing only
magnitude does not change direction. If $g_u=\partial L/\partial u$, then
\[
 \frac{\partial L}{\partial W}=g_u h^\top,
 \qquad \frac{\partial L}{\partial b}=g_u.
\]
Batch gradients add the contributions prescribed by the loss reduction. The
encoder weights do not enter the update. Notebook 1 works through a small
normalized-head gradient numerically; its simple loss is explicitly not the
project's full training objective. This is ordinary backpropagation through
differentiable operations. The hard contextual neighborhood instead uses a
prescribed surrogate gradient, explained in the next section.

\subsection{Why AdamW, and what is fixed?}
The actual optimizer is AdamW. For gradient $g_t$, it keeps exponentially weighted
first and second moments,
\[
 m_t=\beta_1m_{t-1}+(1-\beta_1)g_t,\qquad
 v_t=\beta_2v_{t-1}+(1-\beta_2)g_t^2,
\]
and corrects their startup bias: $\widehat m_t=m_t/(1-\beta_1^t)$ and
$\widehat v_t=v_t/(1-\beta_2^t)$. The essential elementwise update is
\[
 \theta_t=(1-\eta\lambda_w)\theta_{t-1}
       -\eta\frac{\widehat m_t}{\sqrt{\widehat v_t}+\varepsilon_o}.
\]
The first term is decoupled weight decay; the second uses recent gradient
direction and scale. Weight decay is distinct from the contextual recipe's
mean-similarity regularizer, which acts on embeddings through the loss.

All recipes use $\beta_1=0.9$, $\beta_2=0.999$, $\varepsilon_o=10^{-8}$,
and $\lambda_w=10^{-4}$ on weight and bias. Learning rate is constant within
a run. There is no scheduler, dropout, extra augmentation, gradient clipping,
early stopping, or mixed precision. These are fixed controls, not assertions
that these choices are universally optimal.

\studycheck{Explain why a frozen encoder still permits learning, why the head
can recognize an action absent from its training labels, and why perturbing
cached features would change the experiment. Answers: Questions 3--4.}
\source{\src{src/pose_embed/data/motionbert.py}{Deterministic preprocessing};
\src{src/pose_embed/models/motionbert.py}{Encoder adapter};
\src{src/pose_embed/models/action_head.py}{Pooling and head};
\src{src/pose_embed/training/runner.py}{Training runner};
\href{https://github.com/Walter0807/MotionBERT/blob/705d3a95354db8bdb696b3492e47a3b5537174ff/lib/model/DSTformer.py}{Pinned DSTformer source}.}

\section{What is reused, adapted, and written for this study}\label{sec:provenance}
The research contribution is not a new pose detector or a new transformer
backbone. It is the specified transfer comparison and the controlled process
needed to interpret its result.

\begin{longtable}{@{}P{0.21\linewidth}P{0.33\linewidth}P{0.40\linewidth}@{}}
\toprule Component & Provenance & Role and boundary\\\midrule\endhead
Pose data & Existing NTU dataset and HRNet-derived pose annotations distributed
through OpenMMLab/PYSKL & Reuse the available estimates. No new HRNet training,
video collection, or native 25-joint 3D input.\\\addlinespace
MotionBERT encoder & Unchanged code imported from a pinned Apache-2.0 checkout;
verified pretrained checkpoint & Reuse the architecture and learned weights;
neither source clone nor checkpoint is committed into this repository.\\\addlinespace
Preprocessing & Minimal licensed port of MotionBERT action-data utilities & Make
sampling deterministic, preserve confidence through tracking/conversion, and
apply the declared confidence correction.\\\addlinespace
Pooling and head & Ported MotionBERT action-head semantics & Reuse time/person
pooling and projection structure; train fresh head weights for our objectives.\\\addlinespace
Contextual loss & Independently implemented from published equations, developed
from the project's teaching implementation & The authors' source repository is
reference-only because no license was discovered. It is not copied or adapted.\\\addlinespace
Contrastive and SupCon & Local implementations of the declared formulas & Make
reduction rules, margins, temperature, and single-view semantics explicit.\\\addlinespace
Optimization engine & PyTorch tensor operations, autograd, and AdamW & Reuse
numerical machinery; local code defines which parameters and batches it updates.\\\addlinespace
Training and evaluation & Local protocol-specific runners & Train on cached
auxiliary features; keep novel evaluation out of the training loop.\\\addlinespace
Experimental controls & Local manifests, sampler, corruptions, locks, provenance,
analysis, and report code & Define a reproducible experiment rather than inherit
benchmark defaults.\\\bottomrule
\end{longtable}

\subsection{Why not reuse the upstream one-shot training loop?}
That loop evaluates novel classes during training. Our study requires those
outcomes to remain unavailable until training, selection, and the final lock
are complete. Reusing the encoder is compatible with this requirement; reusing
that evaluation schedule is not. The local head-training loop is a necessary
change in scientific procedure, not a claim that upstream code is defective.

\subsection{Two kinds of provenance}
\textbf{Intellectual provenance} identifies whose method, data, and code are
being used, along with the permitted reuse and attribution. \textbf{Execution
provenance} identifies the exact source revision, checkpoint, input IDs,
configuration, and runtime that produced an artifact. A result needs both.

Historical v2 contains other methods and upstream sources, but those are not all
used in this v3 comparison. The current study does not need a tour of every
preserved benchmark implementation. Nor does a source repository being publicly
viewable establish permission to copy its code; the project's source manifest
and notices record that boundary.

\source{\src{third_party/upstreams.toml}{Pinned upstream manifest};
\src{THIRD_PARTY_NOTICES.md}{Third-party notices};
\src{docs/protocol/ntu-input-source-research.md}{Pose-source investigation};
\src{docs/protocol/motionbert-adapter-gate.md}{Adapter reuse boundary};
\src{docs/protocol/contextual-v3-source-notes.md}{Contextual source and equation notes}.}
